\subsection{Common beam crossings with heading retained}

We set the rotating beam to $-90^\circ$ at each scan midpoint and solved for its crossing of the moving target bearing. The four arms shared these body-beam times. At the scan seam the pose delay wrapped to the neighboring midpoint. We tested eight bearings from $-90^\circ$ to $180^\circ$ with and without $1^\circ$ sinusoidal yaw at the wave frequency. Level pointing retained yaw at the assigned pose time.

At zero yaw the original forward counts were recovered with the quarter-scan clock. Solving the crossing changed the forward interaction from $-23.16$ to $-23.19$ points at 0.36125 Hz and 0.376 m. With yaw retained, the stand-alone leveling gain fell from 24.06 to $-1.86$ points, and the interaction became $+2.73$ points. At bearings $(-90,0,45,89.5,90,90.5,135,180)^\circ$, zero-yaw crossing interactions in this wave were $(0,-23.19,-0.47,-7.91,-2.70,-4.25,+13.33,-23.08)$ points. At $135^\circ$ leveling alone lost 13.33 points, reversing the interaction even without yaw. Adding yaw cut the forward body-midpoint usable count from 1,161 to 197 of 6,144 and raised the pose-only gain from 23.19 to 38.43 points. In the reference wave the count fell from 5,712 to 696 and timing alone gained 82.03 points, while the interaction reached $+6.64$ points. Crossing shifts were at most 10.8 ms. The stale heading can therefore govern the benefit of a level mount even when the spoke time changes little.
